% Options for packages loaded elsewhere
\PassOptionsToPackage{unicode}{hyperref}
\PassOptionsToPackage{hyphens}{url}
%
\documentclass[
]{article}
\usepackage{amsmath,amssymb}
\usepackage{iftex}
\ifPDFTeX
  \usepackage[T1]{fontenc}
  \usepackage[utf8]{inputenc}
  \usepackage{textcomp} % provide euro and other symbols
\else % if luatex or xetex
  \usepackage{unicode-math} % this also loads fontspec
  \defaultfontfeatures{Scale=MatchLowercase}
  \defaultfontfeatures[\rmfamily]{Ligatures=TeX,Scale=1}
\fi
\usepackage{lmodern}
\ifPDFTeX\else
  % xetex/luatex font selection
\fi
% Use upquote if available, for straight quotes in verbatim environments
\IfFileExists{upquote.sty}{\usepackage{upquote}}{}
\IfFileExists{microtype.sty}{% use microtype if available
  \usepackage[]{microtype}
  \UseMicrotypeSet[protrusion]{basicmath} % disable protrusion for tt fonts
}{}
\makeatletter
\@ifundefined{KOMAClassName}{% if non-KOMA class
  \IfFileExists{parskip.sty}{%
    \usepackage{parskip}
  }{% else
    \setlength{\parindent}{0pt}
    \setlength{\parskip}{6pt plus 2pt minus 1pt}}
}{% if KOMA class
  \KOMAoptions{parskip=half}}
\makeatother
\usepackage{xcolor}
\usepackage[margin=1in]{geometry}
\usepackage{color}
\usepackage{fancyvrb}
\newcommand{\VerbBar}{|}
\newcommand{\VERB}{\Verb[commandchars=\\\{\}]}
\DefineVerbatimEnvironment{Highlighting}{Verbatim}{commandchars=\\\{\}}
% Add ',fontsize=\small' for more characters per line
\usepackage{framed}
\definecolor{shadecolor}{RGB}{248,248,248}
\newenvironment{Shaded}{\begin{snugshade}}{\end{snugshade}}
\newcommand{\AlertTok}[1]{\textcolor[rgb]{0.94,0.16,0.16}{#1}}
\newcommand{\AnnotationTok}[1]{\textcolor[rgb]{0.56,0.35,0.01}{\textbf{\textit{#1}}}}
\newcommand{\AttributeTok}[1]{\textcolor[rgb]{0.13,0.29,0.53}{#1}}
\newcommand{\BaseNTok}[1]{\textcolor[rgb]{0.00,0.00,0.81}{#1}}
\newcommand{\BuiltInTok}[1]{#1}
\newcommand{\CharTok}[1]{\textcolor[rgb]{0.31,0.60,0.02}{#1}}
\newcommand{\CommentTok}[1]{\textcolor[rgb]{0.56,0.35,0.01}{\textit{#1}}}
\newcommand{\CommentVarTok}[1]{\textcolor[rgb]{0.56,0.35,0.01}{\textbf{\textit{#1}}}}
\newcommand{\ConstantTok}[1]{\textcolor[rgb]{0.56,0.35,0.01}{#1}}
\newcommand{\ControlFlowTok}[1]{\textcolor[rgb]{0.13,0.29,0.53}{\textbf{#1}}}
\newcommand{\DataTypeTok}[1]{\textcolor[rgb]{0.13,0.29,0.53}{#1}}
\newcommand{\DecValTok}[1]{\textcolor[rgb]{0.00,0.00,0.81}{#1}}
\newcommand{\DocumentationTok}[1]{\textcolor[rgb]{0.56,0.35,0.01}{\textbf{\textit{#1}}}}
\newcommand{\ErrorTok}[1]{\textcolor[rgb]{0.64,0.00,0.00}{\textbf{#1}}}
\newcommand{\ExtensionTok}[1]{#1}
\newcommand{\FloatTok}[1]{\textcolor[rgb]{0.00,0.00,0.81}{#1}}
\newcommand{\FunctionTok}[1]{\textcolor[rgb]{0.13,0.29,0.53}{\textbf{#1}}}
\newcommand{\ImportTok}[1]{#1}
\newcommand{\InformationTok}[1]{\textcolor[rgb]{0.56,0.35,0.01}{\textbf{\textit{#1}}}}
\newcommand{\KeywordTok}[1]{\textcolor[rgb]{0.13,0.29,0.53}{\textbf{#1}}}
\newcommand{\NormalTok}[1]{#1}
\newcommand{\OperatorTok}[1]{\textcolor[rgb]{0.81,0.36,0.00}{\textbf{#1}}}
\newcommand{\OtherTok}[1]{\textcolor[rgb]{0.56,0.35,0.01}{#1}}
\newcommand{\PreprocessorTok}[1]{\textcolor[rgb]{0.56,0.35,0.01}{\textit{#1}}}
\newcommand{\RegionMarkerTok}[1]{#1}
\newcommand{\SpecialCharTok}[1]{\textcolor[rgb]{0.81,0.36,0.00}{\textbf{#1}}}
\newcommand{\SpecialStringTok}[1]{\textcolor[rgb]{0.31,0.60,0.02}{#1}}
\newcommand{\StringTok}[1]{\textcolor[rgb]{0.31,0.60,0.02}{#1}}
\newcommand{\VariableTok}[1]{\textcolor[rgb]{0.00,0.00,0.00}{#1}}
\newcommand{\VerbatimStringTok}[1]{\textcolor[rgb]{0.31,0.60,0.02}{#1}}
\newcommand{\WarningTok}[1]{\textcolor[rgb]{0.56,0.35,0.01}{\textbf{\textit{#1}}}}
\usepackage{graphicx}
\makeatletter
\def\maxwidth{\ifdim\Gin@nat@width>\linewidth\linewidth\else\Gin@nat@width\fi}
\def\maxheight{\ifdim\Gin@nat@height>\textheight\textheight\else\Gin@nat@height\fi}
\makeatother
% Scale images if necessary, so that they will not overflow the page
% margins by default, and it is still possible to overwrite the defaults
% using explicit options in \includegraphics[width, height, ...]{}
\setkeys{Gin}{width=\maxwidth,height=\maxheight,keepaspectratio}
% Set default figure placement to htbp
\makeatletter
\def\fps@figure{htbp}
\makeatother
\setlength{\emergencystretch}{3em} % prevent overfull lines
\providecommand{\tightlist}{%
  \setlength{\itemsep}{0pt}\setlength{\parskip}{0pt}}
\setcounter{secnumdepth}{-\maxdimen} % remove section numbering
\ifLuaTeX
  \usepackage{selnolig}  % disable illegal ligatures
\fi
\usepackage{bookmark}
\IfFileExists{xurl.sty}{\usepackage{xurl}}{} % add URL line breaks if available
\urlstyle{same}
\hypersetup{
  pdftitle={Lab Activity},
  pdfauthor={Livio},
  hidelinks,
  pdfcreator={LaTeX via pandoc}}

\title{Lab Activity}
\author{Livio}
\date{2026-02-20}

\begin{document}
\maketitle

\subsection{Activity: Discover how contrasts and centering change your
model}\label{activity-discover-how-contrasts-and-centering-change-your-model}

\textbf{Goal}\\
By fitting models with different encodings, \emph{discover} for
yourself:

\begin{itemize}
\tightlist
\item
  How dummy vs sum contrasts change coefficient \emph{meaning}\\
\item
  Why centering continuous predictors matters (especially with
  interactions)\\
\item
  How to get ``clean'' interpretations of main effects
\end{itemize}

\textbf{Instructions}\\
Work step-by-step. Compare \texttt{summary()} outputs and plots. Explain
what you observe in your own words.

\begin{center}\rule{0.5\linewidth}{0.5pt}\end{center}

\subsubsection{1. Generate the data and visualize
it}\label{generate-the-data-and-visualize-it}

Run:

\begin{Shaded}
\begin{Highlighting}[]
\FunctionTok{set.seed}\NormalTok{(}\DecValTok{1}\NormalTok{)}
\NormalTok{n0 }\OtherTok{\textless{}{-}} \DecValTok{10}
\NormalTok{D }\OtherTok{\textless{}{-}} \FunctionTok{data.frame}\NormalTok{(}
  \AttributeTok{gr =} \FunctionTok{as.factor}\NormalTok{(}\FunctionTok{rep}\NormalTok{(LETTERS[}\DecValTok{1}\SpecialCharTok{:}\DecValTok{2}\NormalTok{], n0)),}
  \AttributeTok{x  =} \FunctionTok{rnorm}\NormalTok{(n0 }\SpecialCharTok{*} \DecValTok{2}\NormalTok{) }\SpecialCharTok{+} \DecValTok{3}
\NormalTok{)}
\NormalTok{mu }\OtherTok{\textless{}{-}} \DecValTok{2} \SpecialCharTok{+}\NormalTok{ (D}\SpecialCharTok{$}\NormalTok{gr }\SpecialCharTok{==} \StringTok{"B"}\NormalTok{) }\SpecialCharTok{*}\NormalTok{ .}\DecValTok{5} \SpecialCharTok{+}\NormalTok{ D}\SpecialCharTok{$}\NormalTok{x }\SpecialCharTok{*}\NormalTok{ (D}\SpecialCharTok{$}\NormalTok{gr }\SpecialCharTok{==} \StringTok{"B"}\NormalTok{) }\SpecialCharTok{*} \DecValTok{2}
\NormalTok{D}\SpecialCharTok{$}\NormalTok{y }\OtherTok{\textless{}{-}}\NormalTok{ mu }\SpecialCharTok{+} \FunctionTok{rnorm}\NormalTok{(n0 }\SpecialCharTok{*} \DecValTok{2}\NormalTok{)}

\FunctionTok{library}\NormalTok{(ggplot2)}
\end{Highlighting}
\end{Shaded}

\begin{verbatim}
## Warning: il pacchetto 'ggplot2' è stato creato con R versione 4.5.3
\end{verbatim}

\begin{Shaded}
\begin{Highlighting}[]
\FunctionTok{ggplot}\NormalTok{(D, }\FunctionTok{aes}\NormalTok{(}\AttributeTok{x =}\NormalTok{ x, }\AttributeTok{y =}\NormalTok{ y, }\AttributeTok{color =}\NormalTok{ gr)) }\SpecialCharTok{+}
  \FunctionTok{geom\_point}\NormalTok{() }\SpecialCharTok{+}
  \FunctionTok{geom\_smooth}\NormalTok{(}\AttributeTok{method =} \StringTok{"lm"}\NormalTok{, }\AttributeTok{fill =} \ConstantTok{NA}\NormalTok{, }\AttributeTok{fullrange =} \ConstantTok{TRUE}\NormalTok{) }\SpecialCharTok{+}
  \FunctionTok{xlim}\NormalTok{(}\SpecialCharTok{{-}}\NormalTok{.}\DecValTok{5}\NormalTok{, }\DecValTok{6}\NormalTok{) }\SpecialCharTok{+}
  \FunctionTok{theme\_bw}\NormalTok{()}
\end{Highlighting}
\end{Shaded}

\begin{verbatim}
## `geom_smooth()` using formula = 'y ~ x'
\end{verbatim}

\includegraphics{student_activity3_files/figure-latex/unnamed-chunk-1-1.pdf}

\textbf{Tasks}

\begin{enumerate}
\def\labelenumi{\arabic{enumi}.}
\tightlist
\item
  Write down (from the code) the \emph{true} generative model for groups
  A and B (two separate regression equations).\\
\item
  From the plot, describe:

  \begin{itemize}
  \tightlist
  \item
    Does \texttt{x} affect \texttt{y} in group A?\\
  \item
    Does \texttt{x} affect \texttt{y} in group B?\\
  \item
    Is there an interaction between \texttt{gr} and \texttt{x}?
  \end{itemize}
\end{enumerate}

\subsubsection{2. Model 1: Default dummy coding (no
centering)}\label{model-1-default-dummy-coding-no-centering}

\begin{Shaded}
\begin{Highlighting}[]
\NormalTok{mod1 }\OtherTok{\textless{}{-}} \FunctionTok{lm}\NormalTok{(y }\SpecialCharTok{\textasciitilde{}}\NormalTok{ gr }\SpecialCharTok{*}\NormalTok{ x, }\AttributeTok{data =}\NormalTok{ D)}
\FunctionTok{summary}\NormalTok{(mod1)}
\end{Highlighting}
\end{Shaded}

\begin{verbatim}
## 
## Call:
## lm(formula = y ~ gr * x, data = D)
## 
## Residuals:
##      Min       1Q   Median       3Q      Max 
## -1.55585 -0.61528 -0.00131  0.54234  1.19203 
## 
## Coefficients:
##             Estimate Std. Error t value Pr(>|t|)   
## (Intercept)  2.30818    1.24368   1.856  0.08198 . 
## grB          1.04292    1.53659   0.679  0.50701   
## x           -0.03137    0.37016  -0.085  0.93352   
## grB:x        1.68703    0.46200   3.652  0.00215 **
## ---
## Signif. codes:  0 '***' 0.001 '**' 0.01 '*' 0.05 '.' 0.1 ' ' 1
## 
## Residual standard error: 0.8778 on 16 degrees of freedom
## Multiple R-squared:  0.9481, Adjusted R-squared:  0.9384 
## F-statistic: 97.49 on 3 and 16 DF,  p-value: 1.708e-10
\end{verbatim}

\begin{Shaded}
\begin{Highlighting}[]
\NormalTok{p2 }\OtherTok{\textless{}{-}} \FunctionTok{ggplot}\NormalTok{(D, }\FunctionTok{aes}\NormalTok{(}\AttributeTok{x =}\NormalTok{ x, }\AttributeTok{y =}\NormalTok{ y, }\AttributeTok{color =}\NormalTok{ gr)) }\SpecialCharTok{+}
  \FunctionTok{geom\_point}\NormalTok{() }\SpecialCharTok{+}
  \FunctionTok{geom\_smooth}\NormalTok{(}\AttributeTok{method =} \StringTok{"lm"}\NormalTok{, }\AttributeTok{fill =} \ConstantTok{NA}\NormalTok{, }\AttributeTok{fullrange =} \ConstantTok{TRUE}\NormalTok{) }\SpecialCharTok{+}
  \FunctionTok{geom\_vline}\NormalTok{(}\AttributeTok{xintercept =} \DecValTok{0}\NormalTok{, }\AttributeTok{linetype =} \StringTok{"dashed"}\NormalTok{) }\SpecialCharTok{+}
  \FunctionTok{xlim}\NormalTok{(}\SpecialCharTok{{-}}\NormalTok{.}\DecValTok{5}\NormalTok{, }\DecValTok{6}\NormalTok{) }\SpecialCharTok{+}
  \FunctionTok{theme\_bw}\NormalTok{()}
\FunctionTok{print}\NormalTok{(p2)}
\end{Highlighting}
\end{Shaded}

\begin{verbatim}
## `geom_smooth()` using formula = 'y ~ x'
\end{verbatim}

\includegraphics{student_activity3_files/figure-latex/unnamed-chunk-2-1.pdf}

\textbf{Discover}\\
1. What does the \texttt{grB} coefficient mean? (Hint: at what
\texttt{x}?)\\
2. Is the group difference meaningful given where data are?\\
3. Look at p-values. Which effects are ``significant''? Interaction?

\begin{center}\rule{0.5\linewidth}{0.5pt}\end{center}

\subsubsection{\texorpdfstring{3. Model 2: Center \texttt{x} (same dummy
coding)}{3. Model 2: Center x (same dummy coding)}}\label{model-2-center-x-same-dummy-coding}

\begin{Shaded}
\begin{Highlighting}[]
\NormalTok{D2 }\OtherTok{\textless{}{-}}\NormalTok{ D}
\NormalTok{D2}\SpecialCharTok{$}\NormalTok{x }\OtherTok{\textless{}{-}}\NormalTok{ D}\SpecialCharTok{$}\NormalTok{x }\SpecialCharTok{{-}} \FunctionTok{mean}\NormalTok{(D}\SpecialCharTok{$}\NormalTok{x)}
\NormalTok{mod2 }\OtherTok{\textless{}{-}} \FunctionTok{lm}\NormalTok{(y }\SpecialCharTok{\textasciitilde{}}\NormalTok{ gr }\SpecialCharTok{*}\NormalTok{ x, }\AttributeTok{data =}\NormalTok{ D2)}
\FunctionTok{summary}\NormalTok{(mod2)}
\end{Highlighting}
\end{Shaded}

\begin{verbatim}
## 
## Call:
## lm(formula = y ~ gr * x, data = D2)
## 
## Residuals:
##      Min       1Q   Median       3Q      Max 
## -1.55585 -0.61528 -0.00131  0.54234  1.19203 
## 
## Coefficients:
##             Estimate Std. Error t value Pr(>|t|)    
## (Intercept)  2.20810    0.27933   7.905 6.47e-07 ***
## grB          6.42543    0.39449  16.288 2.21e-11 ***
## x           -0.03137    0.37016  -0.085  0.93352    
## grB:x        1.68703    0.46200   3.652  0.00215 ** 
## ---
## Signif. codes:  0 '***' 0.001 '**' 0.01 '*' 0.05 '.' 0.1 ' ' 1
## 
## Residual standard error: 0.8778 on 16 degrees of freedom
## Multiple R-squared:  0.9481, Adjusted R-squared:  0.9384 
## F-statistic: 97.49 on 3 and 16 DF,  p-value: 1.708e-10
\end{verbatim}

\begin{Shaded}
\begin{Highlighting}[]
\NormalTok{p3 }\OtherTok{\textless{}{-}} \FunctionTok{ggplot}\NormalTok{(D2, }\FunctionTok{aes}\NormalTok{(}\AttributeTok{x =}\NormalTok{ x, }\AttributeTok{y =}\NormalTok{ y, }\AttributeTok{color =}\NormalTok{ gr)) }\SpecialCharTok{+}
  \FunctionTok{geom\_point}\NormalTok{() }\SpecialCharTok{+}
  \FunctionTok{geom\_smooth}\NormalTok{(}\AttributeTok{method =} \StringTok{"lm"}\NormalTok{, }\AttributeTok{fill =} \ConstantTok{NA}\NormalTok{, }\AttributeTok{fullrange =} \ConstantTok{TRUE}\NormalTok{) }\SpecialCharTok{+}
  \FunctionTok{geom\_vline}\NormalTok{(}\AttributeTok{xintercept =} \DecValTok{0}\NormalTok{, }\AttributeTok{linetype =} \StringTok{"dashed"}\NormalTok{) }\SpecialCharTok{+}
  \FunctionTok{theme\_bw}\NormalTok{()}
\FunctionTok{print}\NormalTok{(p3)}
\end{Highlighting}
\end{Shaded}

\begin{verbatim}
## `geom_smooth()` using formula = 'y ~ x'
\end{verbatim}

\includegraphics{student_activity3_files/figure-latex/unnamed-chunk-3-1.pdf}

\textbf{Discover}\\
1. Now what does \texttt{grB} mean? (Compare to plot.) Better?\\
2. How do p-values change for main effects? For interaction?\\
3. \textbf{Why} did centering \texttt{x} change the \emph{main effect}
tests, even though interaction was always in the model?

\begin{center}\rule{0.5\linewidth}{0.5pt}\end{center}

\subsubsection{\texorpdfstring{4. Model 3: Zero-sum contrasts for
\texttt{gr} (+ centered
\texttt{x})}{4. Model 3: Zero-sum contrasts for gr (+ centered x)}}\label{model-3-zero-sum-contrasts-for-gr-centered-x}

\begin{Shaded}
\begin{Highlighting}[]
\NormalTok{D3 }\OtherTok{\textless{}{-}}\NormalTok{ D2}
\FunctionTok{contrasts}\NormalTok{(D3}\SpecialCharTok{$}\NormalTok{gr) }\OtherTok{\textless{}{-}} \FunctionTok{contr.sum}\NormalTok{(}\DecValTok{2}\NormalTok{)}
\NormalTok{mod3 }\OtherTok{\textless{}{-}} \FunctionTok{lm}\NormalTok{(y }\SpecialCharTok{\textasciitilde{}}\NormalTok{ gr }\SpecialCharTok{*}\NormalTok{ x, }\AttributeTok{data =}\NormalTok{ D3)}
\FunctionTok{summary}\NormalTok{(mod3)}
\end{Highlighting}
\end{Shaded}

\begin{verbatim}
## 
## Call:
## lm(formula = y ~ gr * x, data = D3)
## 
## Residuals:
##      Min       1Q   Median       3Q      Max 
## -1.55585 -0.61528 -0.00131  0.54234  1.19203 
## 
## Coefficients:
##             Estimate Std. Error t value Pr(>|t|)    
## (Intercept)   5.4208     0.1972  27.483 6.80e-15 ***
## gr1          -3.2127     0.1972 -16.288 2.21e-11 ***
## x             0.8121     0.2310   3.516  0.00287 ** 
## gr1:x        -0.8435     0.2310  -3.652  0.00215 ** 
## ---
## Signif. codes:  0 '***' 0.001 '**' 0.01 '*' 0.05 '.' 0.1 ' ' 1
## 
## Residual standard error: 0.8778 on 16 degrees of freedom
## Multiple R-squared:  0.9481, Adjusted R-squared:  0.9384 
## F-statistic: 97.49 on 3 and 16 DF,  p-value: 1.708e-10
\end{verbatim}

\begin{Shaded}
\begin{Highlighting}[]
\NormalTok{p4 }\OtherTok{\textless{}{-}} \FunctionTok{ggplot}\NormalTok{(D3, }\FunctionTok{aes}\NormalTok{(}\AttributeTok{x =}\NormalTok{ x, }\AttributeTok{y =}\NormalTok{ y, }\AttributeTok{color =}\NormalTok{ gr)) }\SpecialCharTok{+}
  \FunctionTok{geom\_point}\NormalTok{() }\SpecialCharTok{+}
  \FunctionTok{geom\_smooth}\NormalTok{(}\AttributeTok{method =} \StringTok{"lm"}\NormalTok{, }\AttributeTok{fill =} \ConstantTok{NA}\NormalTok{, }\AttributeTok{fullrange =} \ConstantTok{TRUE}\NormalTok{) }\SpecialCharTok{+}
  \FunctionTok{geom\_abline}\NormalTok{(}\AttributeTok{intercept =} \FunctionTok{coef}\NormalTok{(mod3)[}\DecValTok{1}\NormalTok{] , }\AttributeTok{slope =} \FunctionTok{coef}\NormalTok{(mod3)[}\DecValTok{3}\NormalTok{], }\AttributeTok{color =} \StringTok{"black"}\NormalTok{, }\AttributeTok{linewidth =} \DecValTok{1}\NormalTok{) }\SpecialCharTok{+}
  \FunctionTok{theme\_bw}\NormalTok{()}
\FunctionTok{print}\NormalTok{(p4)}
\end{Highlighting}
\end{Shaded}

\begin{verbatim}
## `geom_smooth()` using formula = 'y ~ x'
\end{verbatim}

\includegraphics{student_activity3_files/figure-latex/unnamed-chunk-4-1.pdf}

\textbf{Discover}\\
1. What population does the intercept + \texttt{x} slope (black line)
describe?\\
2. How is this different from Models 1--2?\\
3. Compare p-values across all 3 models. Which encoding gives
``cleanest'' tests?

\begin{center}\rule{0.5\linewidth}{0.5pt}\end{center}

\subsubsection{\texorpdfstring{5. Bonus: Rescale \texttt{x} instead of
just
centering}{5. Bonus: Rescale x instead of just centering}}\label{bonus-rescale-x-instead-of-just-centering}

\begin{Shaded}
\begin{Highlighting}[]
\NormalTok{D4 }\OtherTok{\textless{}{-}}\NormalTok{ D}
\NormalTok{D4}\SpecialCharTok{$}\NormalTok{x }\OtherTok{\textless{}{-}} \FunctionTok{scale}\NormalTok{(D}\SpecialCharTok{$}\NormalTok{x)  }\CommentTok{\# Centers AND divides by SD}
\NormalTok{mod4 }\OtherTok{\textless{}{-}} \FunctionTok{lm}\NormalTok{(y }\SpecialCharTok{\textasciitilde{}}\NormalTok{ gr }\SpecialCharTok{*}\NormalTok{ x, }\AttributeTok{data =}\NormalTok{ D4)}
\FunctionTok{summary}\NormalTok{(mod4)}
\end{Highlighting}
\end{Shaded}

\begin{verbatim}
## 
## Call:
## lm(formula = y ~ gr * x, data = D4)
## 
## Residuals:
##      Min       1Q   Median       3Q      Max 
## -1.55585 -0.61528 -0.00131  0.54234  1.19203 
## 
## Coefficients:
##             Estimate Std. Error t value Pr(>|t|)    
## (Intercept)  2.20810    0.27933   7.905 6.47e-07 ***
## grB          6.42543    0.39449  16.288 2.21e-11 ***
## x           -0.02865    0.33805  -0.085  0.93352    
## grB:x        1.54069    0.42193   3.652  0.00215 ** 
## ---
## Signif. codes:  0 '***' 0.001 '**' 0.01 '*' 0.05 '.' 0.1 ' ' 1
## 
## Residual standard error: 0.8778 on 16 degrees of freedom
## Multiple R-squared:  0.9481, Adjusted R-squared:  0.9384 
## F-statistic: 97.49 on 3 and 16 DF,  p-value: 1.708e-10
\end{verbatim}

\textbf{Discover}\\
1. Compare \texttt{mod4} coefficients to \texttt{mod2} (just centering).
Why are slopes \textasciitilde2x larger now?\\
2. Is the group difference (\texttt{grB}) at \texttt{x=0} the
\emph{same}?\\
3. \textbf{Pro}: What does ``per 1-SD change in \texttt{x}'' mean for
interpretation?\\
4. \textbf{Con}: Does rescaling change p-values or tests? Why/why not?
When would you rescale?

\begin{center}\rule{0.5\linewidth}{0.5pt}\end{center}

\subsubsection{6. Your discoveries}\label{your-discoveries}

\textbf{Answer these (no peeking at solutions)}:\\
1. Why is the ``main effect'' of \texttt{gr} meaningless without
specifying \emph{at what \texttt{x}}?\\
2. How does centering \texttt{x} make main effects interpretable
\textbf{despite} interaction?\\
3. When interactions are present, why do defaults hurt
power/interpretation?\\
4. Rule of thumb: when to center? When to use \texttt{contr.sum()}?

\textbf{Extension (groups)}: Fit on \emph{your own} simulated data
(change slopes/groups). Does the lesson hold?

\begin{center}\rule{0.5\linewidth}{0.5pt}\end{center}

\subsection{Answer to your side question: Rescaling vs
centering}\label{answer-to-your-side-question-rescaling-vs-centering}

Rescaling (\texttt{scale(x)} = centering + dividing by SD) has these
effects:

\textbf{Pros}:\\
- Coefficients are \textbf{standardized effect sizes} (``per 1-SD change
in \texttt{x}''). Great for comparing predictor strength (e.g., BMI vs
age). - Group difference still at meaningful \texttt{x=0} (mean).

\textbf{Cons vs pure centering}:\\
- \textbf{No change to p-values, F-tests, or R²} (pure linear
rescaling). Tests identical to centered-only model.\\
- Slopes scale up by \(1/\text{SD}_x\), so intercept adjusts down
accordingly.\\
- Less intuitive if SD lacks meaning (e.g., arbitrary imaging
parameter).

\textbf{When to rescale}: For predictor comparison or publication
tables. Otherwise, centering suffices for clean main effects.

\end{document}
